  \subsubsection{整数性と不等式で整数部分を確定する}\label{節：ガウス型の整数部分の確定}
    \bookterm{floor}{ガウス記号}を見たら,まず中にある数が整数かどうかを考えよう.
    整数$k$と実数$x$に対して,
    \[
      \lfloor k+x\rfloor=k+\lfloor x\rfloor
    \]
    が成り立つ.例えば,整数$a_n$に対して$\lfloor a_n+\sqrt3\rfloor=a_n+1$である.
    ただし,これを使う前に,初項から整数が保たれることを確かめる必要がある.

    中身が整数との和でなくても,\,$k\le X<k+1$を示せば$\lfloor X\rfloor=k$と分かる.
    平方根なら,先に隣り合う平方数の間に挟むとよい.
    \textbf{\bookterm{floor-function}{床関数}の値を確定してから,既知の\bookterm{recurrence}{漸化式}として解く}のがこの使い方である.

    \noindent \bookblocklink{example-gauss-role-1}{problem}{answer}{【例題】}\par\nobreak
    \begin{enumerate}[label=(\arabic*),parsep=3pt,beginpenalty=10000,format=\bookitemlink{example-gauss-role-1}{problem}{answer}]
      \item \bookterm{sequence}{数列}を次のように定める.\bookterm{general}{一般項}を求めよ.
\BookMathLayout{            \[
              a_1=1,\qquad
              a_{n+1}=\left\lfloor\sqrt{a_n^2+4a_n+5}\right\rfloor
              \quad(n\ge1).
            \]}{            \[
              \begin{aligned}
              a_1&=1,\\
              a_{n+1}&=\left\lfloor\sqrt{a_n^2+4a_n+5}\right\rfloor
              \\              &\quad(n\ge1).
              \end{aligned}
            \]}
    \end{enumerate}

    \noindent \bookblocklink{example-gauss-role-1}{hint}{problem}{【方針】}\par\nobreak
    \begin{enumerate}[label=(\arabic*),parsep=3pt,beginpenalty=10000,format=\bookitemlink{example-gauss-role-1}{hint}{problem}]
      \item 整数$a_n$に対して,平方根の中身を$(a_n+2)^2$と$(a_n+3)^2$で挟みます.
            まず各項が非負整数であることを確かめましょう.
    \end{enumerate}

    \noindent \bookblocklink{example-gauss-role-1}{answer}{problem}{【解答】}\par\nobreak
    \begin{enumerate}[label=(\arabic*),parsep=3pt,beginpenalty=10000,format=\bookitemlink{example-gauss-role-1}{answer}{problem}]
      \item \[
              a_n=2n-1\quad(n\ge1).
            \]
    \end{enumerate}

    \noindent \bookblocklink{example-gauss-role-1}{solution}{problem}{【解法】}\par\nobreak
    \begin{enumerate}[label=(\arabic*),parsep=3pt,beginpenalty=10000,format=\bookitemlink{example-gauss-role-1}{solution}{problem}]
      \item 初項は非負整数であり,非負整数$a_n$から次の項も非負整数になる.
            よって帰納法により,すべての項が非負整数である.
            非負整数$x$について,
            \[
              (x+2)^2\le x^2+4x+5<(x+3)^2
            \]
            が成り立つ.左の差は1,右の差は$2x+4>0$である.
            平方根を取ると,
            \[
              x+2\le\sqrt{x^2+4x+5}<x+3
            \]
            なので,\bookterm{floor-function}{床関数}の値は$x+2$である.
            従って$a_{n+1}=a_n+2$となり,\sectionref{節：等差型}に帰着して$a_n=2n-1$を得る.
    \end{enumerate}
